\chapter{Introducción}
\label{cap:introduccion}

En los programas de posgrado universitario, la elaboración de la tesis constituye una de las etapas académicas más determinantes para la formación de los estudiantes, donde la comunicación continua y el acompañamiento con el asesor son fundamentales para evitar que los alumnos se atrasen. Sin embargo, este proceso suele verse afectado por el uso de herramientas informales de comunicación, tales como correos electrónicos personales y aplicaciones de mensajería. Esta práctica habitual provoca la pérdida de archivos, confusiones entre las distintas versiones de un mismo documento y el olvido frecuente de los acuerdos tomados en las asesorías debido a la falta de minutas organizadas. A su vez, los sistemas informáticos que ya tiene la universidad se concentran casi en su totalidad en trámites administrativos generales (como inscripciones y control escolar), sin contar con apartados específicos para dar seguimiento al trabajo continuo de la tesis ni avisar a tiempo a las coordinaciones cuando un alumno pasa mucho tiempo inactivo.

Para dar una respuesta a esta problemática, en este trabajo se propone el diseño y desarrollo de una página web orientada a reunir y organizar el seguimiento de las actividades de tesis. Con el fin de asegurar que el proyecto se termine dentro de los tiempos fijados para la tesina, el alcance de esta primera versión funcional se delimita a tres funciones esenciales: Recepción y consulta de entregables en formato PDF, Registro digital de minutas, Supervisión de fechas y control de inactividad. 

% =========================================================================
% LLAMADA A LAS SECCIONES HIJAS CON SALTO DE PÁGINA
% =========================================================================

\clearpage
\section{Resumen}\label{sec:resumen}El presente proyecto consiste en el diseño y desarrollo de una \textbf{página web} orientada a facilitar el seguimiento, control y organización de las actividades de tesis en programas de posgrado. El alcance de esta primera versión del sistema se delimita a tres funciones esenciales: la recepción y consulta de avances en formato PDF guardando el historial de cambios sin sobreescribir archivos; el registro digital de minutas al término de cada asesoría con los acuerdos y tareas establecidos; y la revisión básica de fechas para identificar y avisar oportunamente cuando un alumno sume más de 30 días de inactividad académica. La propuesta reúne estas herramientas en un solo lugar de manera accesible y segura para alumnos, asesores y coordinadores, sentando las bases operativas del acompañamiento académico para evitar atrasos en la titulación de los tesistas.  \vspace{0.5cm}

\noindent \textbf{Palabras clave:} seguimiento de tesis, apoyo escolar, control de archivos, minutas de asesoría, página web educativa, prevención del rezago
\clearpage
\section{Objetivos del Proyecto}
\label{sec:objetivos}

\subsection{Objetivo General}
\label{subsec:obj_general}
    Desarrollar una primera versión de una \textbf{página web} para organizar, centralizar y dar seguimiento continuo a las actividades de tesis en los programas de posgrado, reuniendo en un solo lugar las herramientas necesarias para que alumnos, asesores y coordinadores trabajen de forma coordinada. 

    Para lograr esto, la página web se enfocará en resolver tres necesidades principales:

    \textbf{Control de documentos:} Permitir que los estudiantes suban sus borradores de tesis exclusivamente en formato PDF, guardando un historial ordenado de cada versión entregada para que el asesor pueda revisarlas sin riesgo de sobreescribir correcciones anteriores ni borrar archivos por error.

    \textbf{Registro de asesorías:} Facilitar el llenado y la consulta de minutas digitales al terminar cada sesión de trabajo, dejando por escrito la fecha, los temas tratados y la lista de acuerdos o tareas pendientes que deben cumplir tanto el alumno como el profesor.

    \textbf{Supervisión y prevención del rezago:} Llevar un control automático de las fechas para detectar si un estudiante acumula más de 30 días sin registrar avances o reuniones, mostrando avisos claros para que la coordinación y los tutores puedan intervenir a tiempo y evitar que los alumnos se atrasen en su titulación.
\subsection{Objetivos Específicos}
\label{subsec:obj_especificos}

Para cumplir con el objetivo general, el proyecto se divide en las siguientes etapas y actividades específicas:

\begin{enumerate}
    \item \textbf{Diseño de accesos y pantallas según el tipo de usuario (Roles):}
    \begin{itemize}
        \item Diseñar pantallas sencillas, claras y fáciles de usar para que cualquier persona pueda navegar sin perderse.
        \item Separar el acceso en tres tipos de cuentas: el estudiante (para subir sus trabajos y revisar tareas), el asesor (para dejar comentarios y calificar acuerdos) y la coordinación (para ver el estado general de todos los alumnos).
        \item \textit{Restricción:} Por seguridad y privacidad, cada usuario solo podrá ver su propia información; ningún estudiante podrá ver los archivos o minutas de otro compañero.
    \end{itemize}

    \item \textbf{Módulo de entrega y control de documentos de tesis:}
    \begin{itemize}
        \item Permitir que el alumno suba sus entregas a la plataforma para que el profesor las pueda consultar y descargar cuando las vaya a revisar.
        \item Llevar un registro ordenado que guarde cada versión del documento con su fecha y hora, mostrando un historial claro de los cambios realizados.
        \item \textit{Restricciones:} 
        \begin{itemize}
            \item El sistema solo aceptará archivos en formato PDF para asegurar que el contenido no se desacomode ni se altere.
            \item Queda prohibido sobreescribir o borrar versiones anteriores; cada nuevo documento se guardará como una versión adicional (por ejemplo: Versión 1, Versión 2) para conservar el historial completo.
        \end{itemize}
    \end{itemize}

    \item \textbf{Módulo para el registro y consulta de minutas de asesoría:}
    \begin{itemize}
        \item Crear un formulario digital donde el asesor o el alumno puedan registrar rápidamente lo sucedido al terminar cada sesión de trabajo.
        \item Dejar por escrito los temas que se revisaron, la fecha de la reunión y la lista de acuerdos o tareas pendientes para la siguiente cita.
        \item \textit{Restricciones:} Toda minuta debe incluir obligatoriamente la fecha de la sesión y al menos un acuerdo o compromiso registrado para poder ser guardada en la página web.
    \end{itemize}

    \item \textbf{Módulo de fechas y avisos de inactividad para evitar atrasos:}
    \begin{itemize}
        \item Monitorear las fechas del sistema calculando de forma automática los días transcurridos desde la última entrega de borrador o desde la última minuta registrada.
        \item Mostrar avisos visuales claros a los tutores y a la coordinación cuando un estudiante lleve mucho tiempo sin interactuar en el sistema.
        \item \textit{Restricción:} El sistema marcará un aviso de inactividad únicamente cuando el alumno acumule más de 30 días consecutivos sin subir un avance ni registrar una minuta de asesoría.
    \end{itemize}

    \item \textbf{Pruebas de funcionamiento del sistema:}
    \begin{itemize}
        \item Probar individualmente cada parte de la página web (subida de archivos, llenado de minutas y cálculo de días) para verificar que funcionen sin errores.
        \item Confirmar con ejemplos reales que los archivos se descarguen bien, que las cuentas respeten sus permisos y que las alertas se activen en las fechas correctas.
    \end{itemize}
\end{enumerate}
\section{Justificación}
\label{sec:justificacion}

La elaboración de una tesis de posgrado requiere una comunicación constante entre el estudiante y su asesor. Sin embargo, en la práctica diaria este proceso suele llevarse a cabo por medios informales como aplicaciones de mensajería instantánea y correos electrónicos personales. Aunque estas herramientas son de uso común, no fueron diseñadas para la gestión académica. Como resultado, los borradores se mezclan entre tantos mensajes, se borran archivos por error, no se sabe con certeza cuál es la versión más reciente del trabajo y se olvidan las tareas acordadas porque no existe un registro formal de las reuniones.

A su vez, los sistemas informáticos con los que cuenta la institución se enfocan casi por completo en trámites escolares generales (como inscripciones y calificaciones), dejando sin apoyo el seguimiento continuo del trabajo de investigación. Por esta razón, se justifica el diseño y desarrollo de una \textbf{página web} que concentre en un solo lugar digital todo el proceso de asesorías, resolviendo de manera directa las necesidades de cada participante:

\begin{enumerate}
    \item \textbf{Beneficio para los estudiantes:} Les proporciona un espacio seguro donde subir sus avances de tesis en formato PDF, con la tranquilidad de que sus entregas anteriores quedan guardadas en un historial y no se van a perder ni sobreescribir. Además, les permite revisar en cualquier momento las minutas con los compromisos y fechas límite acordados con su profesor.
    
    \item \textbf{Beneficio para los asesores:} Les facilita revisar los documentos ordenadamente por versiones, retroalimentar el trabajo y registrar en pocos minutos los acuerdos al término de cada sesión, evitando confusiones y malos entendidos sobre las tareas pendientes.
    
    \item \textbf{Beneficio para la coordinación de posgrado:} Le permite saber el estado real del avance de los estudiantes sin tener que pedir reportes manuales a cada maestro. Gracias al control de fechas, la coordinación puede identificar de inmediato qué alumnos tienen más de 30 días de inactividad, pudiendo intervenir a tiempo para evitar que abandonen la tesis o se atrasen en su titulación.
\end{enumerate}

Por último, delimitar el alcance a una primera versión centrada en tres funciones básicas (control de archivos PDF sin sobreescritura, minutas digitales y alertas por inactividad mayor a 30 días) garantiza la viabilidad del proyecto. Esto permite concluir el sistema dentro de los tiempos estipulados para la tesina, entregando una herramienta funcional, fácil de utilizar y útil para la facultad, sobre la cual se podrán agregar más funciones en desarrollos futuros.